\section{Lifting and the cost of reusing a direction}\label{sec:lifting}

Let $C\subset S^{D-1}$ be a kissing configuration. Decompose the enlarged
space as $\R^{D+k}=\R^D\oplus\R^k$. We refer to the second component as
the \emph{tail}. Fix $\ell\ge1$ distinct tails $t_1,\ldots,t_\ell$ with
$0<\norm{t_i}<1$, and put $a_i=\sqrt{1-\norm{t_i}^2}$. Choose subsets
$A_i\subseteq C$, write $U=\bigcup_i A_i$, and let
$Y\subset S^{k-1}$ be a kissing configuration. Consider
\begin{align}
 Z&=\{(x,0):x\in C\setminus U\},\nonumber\\
 M&=\bigcup_{i=1}^{\ell}\{(a_i x,t_i):x\in A_i\},\label{eq:general-lift}\\
 P&=\{(0,y):y\in Y\}.\nonumber
\end{align}
Put $X=Z\cup M\cup P$. First-component norms distinguish the three
parts, and tails distinguish the copies in $M$, so
\begin{equation}\label{eq:lift-count}
 |Z\cup M\cup P|=|C|-|U|+\sum_i|A_i|+|Y|.
\end{equation}

\begin{theorem}[Labelled lifting criterion]\label{lem:lift}
The set in~\eqref{eq:general-lift} is a kissing configuration if and only if
\begin{align}
 a_i a_j\ip{x}{x'}+\ip{t_i}{t_j}&\le\frac12
 &&\text{whenever }(x,i)\ne(x',j),\ x\in A_i,\ x'\in A_j,
 \label{eq:lift-mixed}\\
 \ip{t_i}{y}&\le\frac12
 &&(i\text{ with }A_i\ne\varnothing,\ y\in Y).\label{eq:lift-pure}
\end{align}
\end{theorem}
\begin{proof}
Every point has norm one. Equator--equator and pure-tail pairs inherit
the inequalities for $C$ and $Y$. Equator--pure-tail pairs are orthogonal.
For an equator--mixed pair, the two original directions are different
because the equator excludes $U$; their inner product is at most
$a_i/2<1/2$. The remaining pairs have exactly the products in
\eqref{eq:lift-mixed} and~\eqref{eq:lift-pure}, which are therefore
both necessary and sufficient.
\end{proof}

Equivalently, the tail Gram matrix specifies a threshold matrix
\begin{equation}\label{eq:kappa}
 \kappa_{ij}=\frac{1/2-\ip{t_i}{t_j}}{a_i a_j},
\end{equation}
and the selected base directions must satisfy
$\ip{x}{x'}\le\kappa_{ij}$ for every distinct labelled pair.
This allows unequal latitudes, subject separately to~\eqref{eq:lift-pure}.

For equal tail norms, put $\norm{t_i}^2=h^2$ and
$a_i=a=\sqrt{1-h^2}$. The same-label condition is
\begin{equation}\label{eq:block-threshold}
 \ip{x}{x'}\le\frac{1/2-h^2}{1-h^2}\qquad(x\ne x',\ x,x'\in A_i),
\end{equation}
whereas two copies of the same direction require
\begin{equation}\label{eq:reuse-threshold}
 \ip{t_i}{t_j}\le h^2-\frac12\qquad(x\in A_i\cap A_j,\ i\ne j).
\end{equation}
These inequalities control block size and multiplicity, respectively.

\subsection{The incidence-capacity inequality}
For general tail norms define a graph $\Gamma$ on the labels by
\begin{equation}\label{eq:reuse-graph}
 ij\in E(\Gamma)\quad\Longleftrightarrow\quad
 a_i a_j+\ip{t_i}{t_j}\le\frac12\qquad(i\ne j).
\end{equation}
The labels assigned to one base direction form a clique in $\Gamma$.

\begin{proposition}[Weighted incidence bound]\label{prop:weighted}
Suppose~\eqref{eq:general-lift} is a kissing configuration with
$|A_i|\le\alpha_i$. If nonnegative numbers $\lambda_1,\ldots,\lambda_\ell$
satisfy
\begin{equation}\label{eq:clique-cover}
 \sum_{i\in J}\lambda_i\ge |J|-1
 \quad\text{for every nonempty clique }J\text{ of }\Gamma,
\end{equation}
then
\begin{equation}\label{eq:weighted-capacity}
 |X|\le |C|+|Y|+\sum_i\lambda_i\alpha_i.
\end{equation}
Equality holds exactly when every used direction's label clique is
tight in~\eqref{eq:clique-cover} and every label with
$\lambda_i>0$ reaches its capacity.
\end{proposition}
\begin{proof}
For $x\in U$, let $J(x)=\{i:x\in A_i\}$. The count gives
\[
 |X|-|C|-|Y|=\sum_{x\in U}(|J(x)|-1)
 \le\sum_{x\in U}\sum_{i\in J(x)}\lambda_i
 =\sum_i\lambda_i|A_i|\le\sum_i\lambda_i\alpha_i.
\]
Both inequalities are sums of nonnegative deficits. Their equality
conditions are precisely those stated.
\end{proof}

If $\Gamma$ has clique number $\omega$, the uniform choice
$\lambda_i=1-1/\omega$ is feasible. Thus
\begin{equation}\label{eq:clique-capacity}
 |X|\le |C|+|Y|+
 \left(1-\frac1\omega\right)\sum_i\alpha_i.
\end{equation}
Minimizing $\sum_i\lambda_i\alpha_i$ subject to
\eqref{eq:clique-cover} is a finite linear program; any rational feasible
choice bounds the prescribed family. The twelve tails below have eight
maximum cliques, all triangles; four disjoint triangles saturate the
uniform bound. Rankin's simplex argument~\cite{rankin} also bounds
multiplicity directly from the tail geometry.

\begin{proposition}[Multiplicity bound]\label{prop:multiplicity}
Suppose~\eqref{eq:general-lift} is a kissing configuration, all tails
have squared norm $h^2$, and $h^2<1/2$. The number of copies of any
original direction is at most
\begin{equation}\label{eq:max-multiplicity}
 r=\min\left\{k+1,\left\lfloor\frac1{1-2h^2}\right\rfloor\right\}.
\end{equation}
If $|A_i|\le\alpha_i$, its cardinality is consequently at most
\begin{equation}\label{eq:reuse-bound}
 |C|+|Y|+\left(1-\frac1r\right)\sum_i\alpha_i.
\end{equation}
\end{proposition}
\begin{proof}
For any clique of $m$ labels,~\eqref{eq:reuse-threshold} gives
\[
 0\le\norm{\sum_{i=1}^{m}t_i}^2
 \le mh^2+m(m-1)(h^2-1/2)
 =\frac m2\bigl(1-m(1-2h^2)\bigr).
\]
This proves the latitude bound. The tails have strictly negative mutual products.
They are affinely independent: an affine dependence would equate
nonempty positive combinations of disjoint sets of tails, whose mutual
inner product is negative, although they are the same vector.
Thus $m\le k+1$ and $\omega(\Gamma)\le r$.
Equation~\eqref{eq:clique-capacity} proves~\eqref{eq:reuse-bound}.
\end{proof}

The multiplicity bound is sharp for one direction at a prescribed
height. For $r\ge2$, a regular simplex of $r$ tails of squared norm
$h^2$ fits in $\R^k$ and has mutual products
$-h^2/(r-1)\le h^2-1/2$; for $r=1$, one tail suffices.
This does not assert simultaneous saturation of block and pure-tail
conditions.

\begin{proposition}[Sharp simplex specialization]\label{prop:simplex}
Let $r\ge2$, let $C\subset S^{D-1}$ be a kissing configuration, and
let $B\subset C$ be nonempty. An equal-latitude lifting replacing
every direction of $B$ by copies at the same $r$ tails, with
$0<h^2<1/2$, exists in $\R^{D+k}$ if and only if
\[
 k\ge r-1,\qquad
 \ip{x}{x'}\le\frac1{r+1}\quad(x\ne x',\ x,x'\in B).
\]
No pure-tail points are required. In particular, there is a kissing
configuration in $\R^{D+r-1}$ with
$|C|+(r-1)|B|$ points. It replaces each direction of $B$ by $r$
copies with tails of squared norm $(r-1)/(2r)$.
\end{proposition}
\begin{proof}
Necessity follows from Proposition~\ref{prop:multiplicity}:
$k\ge r-1$ and $h^2\ge(r-1)/(2r)$. The same-label threshold
$(1/2-h^2)/(1-h^2)$ decreases with $h^2$, so is at most $1/(r+1)$.
For sufficiency, in the hyperplane $\sum_i z_i=0$ of $\R^r$, put
\[
 t_i=\frac1{\sqrt2}\left(e_i-\frac1r\boldsymbol1\right),
 \qquad a=\sqrt{\frac{r+1}{2r}}\qquad(1\le i\le r).
\]
The tails satisfy $\norm{t_i}^2=(r-1)/(2r)$ and
$\ip{t_i}{t_j}=-1/(2r)$ for $i\ne j$. Retain $C\setminus B$
at the equator and use all $(ax,t_i)$ with $x\in B$.
Same-label, distinct-direction products are at most
\[
 \frac{r+1}{2r}\frac1{r+1}+\frac{r-1}{2r}=\frac12.
\]
Distinct labels on the same direction give
$(r+1)/(2r)-1/(2r)=1/2$, and distinct labels on different
directions give at most zero. Theorem~\ref{lem:lift} proves
validity, and~\eqref{eq:lift-count} gives the count.
\end{proof}

At this tail norm the multiplicity bound is $r$, and equality in
its norm-of-sum proof forces the used tails to sum to zero, with
every mutual inner product $-1/(2r)$. Their Gram matrix has rank
$r-1$, so this equality geometry is precisely a regular simplex.
At $h^2=1/4$ the block threshold is $1/3$ and a direction can occur
twice. At $h^2=1/3$ the threshold is $1/4$ and a direction can occur
three times. The Leech constructions realize both cases, fitting
several triangles and pure tails into three tail dimensions in the latter.
The former extends to disjoint blocks as follows.

\begin{theorem}[Antipodal block lift]\label{thm:block-lift}
Let $C\subset S^{D-1}$ be a kissing configuration, and let
$B_1,\ldots,B_t\subset C$ be pairwise disjoint subsets satisfying
\[
 \ip{x}{x'}\le\frac13\qquad(x\ne x',\ x,x'\in B_i).
\]
Let $Y\subset S^{k-1}$ be a kissing configuration containing $t$
distinct antipodal pairs $\{\pm u_i\}$. Then
\begin{align}
 X={}&\{(x,0):x\in C\setminus\textstyle\bigcup_iB_i\}\nonumber\\
 &\quad\cup\{(\tfrac{\sqrt3}{2}x,\tfrac\eps2u_i):
                     x\in B_i,\ \eps\in\{\pm1\}\}
       \cup\{(0,y):y\in Y\}\label{eq:antipodal-lift}
\end{align}
is a kissing configuration in $\R^{D+k}$ with
\begin{equation}\label{eq:antipodal-count}
 |X|=|C|+\sum_i|B_i|+|Y|.
\end{equation}
\end{theorem}
\begin{proof}
Use Theorem~\ref{lem:lift} with tails $\pm u_i/2$. A mixed point has
inner product at most $1/2$ with every pure tail by Cauchy--Schwarz.
For two mixed points in one block at the same tail, the bound is
$(3/4)(1/3)+1/4=1/2$. At opposite tails, different original directions
give at most zero, and equal directions give $3/4-1/4=1/2$.

For points from different blocks, the original directions are distinct,
as are the directed tail labels. Their inner product is at most
$(3/4)(1/2)+(1/4)(1/2)=1/2$. Thus all conditions of the theorem hold.
Each selected direction is deleted once and used twice, giving the count.
\end{proof}

For one block and $Y=\{-1,1\}$ this construction lifts to two latitudes
and includes both poles. In fact their presence fixes the height.
If $(\sqrt{1-h^2}x,\pm h)$ are both used, their mutual inner product
$1-2h^2$ requires $h^2\ge1/4$. The poles require $|h|\le1/2$.
Thus $|h|=1/2$, for $0<|h|<1$. The same-latitude inequality at this
height is exactly the block threshold $1/3$ in the theorem.

Each selected direction adds one point to the orthogonal union of $C$
and $Y$. Section~\ref{app:p48} applies this to large disjoint blocks;
Section~\ref{sec:motion25} varies the two-latitude geometry; and
Section~\ref{sec:leech} uses triple reuse, adding two points per direction.
